% ==================== 第 9 章 总结与展望 ====================
\section{总结与展望}

\subsection{第一阶段成果总结}

本项目完成了 MVP 第一阶段的核心目标，
主要成果如下：

\begin{enumerate}[leftmargin=2em,itemsep=0.2em]
\item \textbf{核心算法：}基于 OR-Tools CP-SAT 9.15 实现排课求解器，
  支持均衡、学生体验、教师稳定、校区效率和局部重排 5 类策略；
  可行方案必须满足全部硬约束，并由独立校验器复核；
\item \textbf{软目标拆解：}将偏好、容量、紧凑度、稳定性、变更数
  六类软目标形式化为可比较的 CP-SAT 线性目标，
  实现了"硬约束严格 + 软目标可解释"；
\item \textbf{稳定性线性化：}提出用辅助变量实现教师日均方差
  的一阶泰勒线性近似，丰富了 CP-SAT 在"分组均衡"
  软目标上的应用；
\item \textbf{最小影响调课：}实现了"教师请假"场景下的最小
  影响重排算法，变更数 $\le 4$；
\item \textbf{REST API：}提供 6 个端点（健康/三方案/仪表盘/
  自定义求解/演示数据/调课），前端与算法解耦；
\item \textbf{前端 4 页面：}决策驾驶舱、智能排课、调课中心、
  资源概览；
\item \textbf{测试体系：}19 个自动化测试（其中 7 个为 API 集成测试）+
  GitHub Actions 双版本 CI 流水线；
\item \textbf{工程化：}一键启动 / 停止 / 跨平台脚本，
  完整中文文档。
\end{enumerate}

\subsection{技术创新点}

本项目在以下三方面具有创新性：

\begin{enumerate}[leftmargin=2em,itemsep=0.2em]
\item \textbf{多目标排课的统一框架：}通过策略权重表把六类软目标
  统一到 CP-SAT 线性目标函数中，使同一份业务数据可生成
  多套可比较的候选方案，为教务决策提供量化依据。
\item \textbf{稳定性软目标的线性化：}把二次目标"教师日均方差"
  转化为一阶泰勒近似的线性目标，保持 CP-SAT 的可解性，
  同时具备非线性的偏好。
\item \textbf{最小影响调课的求解器实现：}复用求解器，仅切换权重
  即可"零成本"支持重排，相比传统"局部搜索"实现更简洁。
\end{enumerate}

\subsection{应用价值}

本项目在 K12 一对一辅导行业具有明确的应用价值：
\begin{itemize}[leftmargin=2em,itemsep=0.1em]
\item \textbf{降本潜力：}减少人工枚举与重复核对工作；
\item \textbf{提质：}演示数据中 8 个课次完整排入且硬冲突数为 0；
\item \textbf{增效：}自动输出调课方案、受影响班级和变更明细；
\item \textbf{透明：}软目标拆解让教务能看到"为什么这样排"，
  减少家校沟通矛盾。
\end{itemize}

\subsection{存在的不足}

本项目目前还存在以下不足：
\begin{enumerate}[leftmargin=2em,itemsep=0.2em]
\item \textbf{业务规模受限：}目前演示数据 8 课次，工业级 K12 机构
  单周可达 1000+ 课次，需要进一步做性能优化与分布式部署。
\item \textbf{权限控制缺失：}所有 API 端点无身份认证，仅适合本地
  试用或内网部署。
\item \textbf{自然语言解析未实现：}目前仅支持 JSON 数据导入，
  不支持"周一 9 点张老师高一数学课排一下"自然语言直接排课。
\item \textbf{飞书生态集成未完成：}第一阶段仅完成了"算法核心 +
  本地驾驶舱"，第二阶段需要继续完成多维表格同步、
  消息卡片、审批、AI 助手等集成。
\end{enumerate}

\subsection{第二阶段展望}

第二阶段（2026.08--2026.12）将重点完成以下工作：
\begin{enumerate}[leftmargin=2em,itemsep=0.2em]
\item \textbf{飞书多维表格双向同步：}业务数据可编辑、课表可回写；
\item \textbf{飞书消息卡片推送：}方案确认后自动 @ 师生；
\item \textbf{飞书审批流程：}调课需教务主任审批后方可发布；
\item \textbf{自然语言解析：}接入飞书 AI 助手，
  教务可用自然语言触发排课。
\end{enumerate}

完成第二阶段后，本项目将形成面向飞书生态的智能排课应用原型，
并在真实业务数据验证、权限与合规机制完善后评估规模化部署可行性。

\subsection{个人收获}

通过本项目的研发，团队成员在以下方面得到锻炼：
\begin{itemize}[leftmargin=2em,itemsep=0.1em]
\item 深入理解约束规划（CP-SAT）的工程化应用；
\item 掌握 Python 数据类、标准库 HTTP、原生前端的现代工程实践；
\item 提升 CI/CD 设计与问题分解能力；
\item 培养产品思维：从"算法能跑"到"用户能懂"的转变。
\end{itemize}

\subsection{本章小结}

本章总结了第一阶段的成果、技术创新、应用价值、
不足与第二阶段规划。本项目已完成 MVP 第一阶段，
具备参赛条件与初步商业化基础，期待在第二阶段
与飞书生态深度融合。
